% Source: https://artofproblemsolving.com/community/c3h1966922_geometry_q
% Stable identity: aops-topic-1966922-q01
The frame is the set of points outside the triangle at distance at most $2$ cm from it. It consists of
\begin{itemize}
\item three rectangles of width $2$ cm along the sides, with total area $2\cdot(\text{perimeter})$, and
\item three circular sectors of radius $2$ cm at the corners. The angle of the sector at a vertex is $180^\circ$ minus the interior angle, so the three sector angles add up to $3\cdot180^\circ-180^\circ=360^\circ$: together they form one full circle of radius $2$.
\end{itemize}

\begin{asy}
usepackage("amsmath");
size(6cm);
defaultpen(fontsize(10pt));
pair B=(0,0),C=(7,0),A=(29/7,sqrt(25-(29/7)^2));
real d=0.4;
pair[] P={A,B,C};
filldraw(A--B--C--cycle,gray(0.8),linewidth(1));
for(int i=0;i<3;++i){
  pair X=P[i],Y=P[(i+1)%3];
  pair n=unit(rotate(-90)*(Y-X));
  draw((X+d*n)--(Y+d*n),linewidth(0.8));
  draw(X--(X+d*n),dotted);draw(Y--(Y+d*n),dotted);
}
for(int i=0;i<3;++i){
  pair X=P[i],Wp=P[(i+2)%3],Y=P[(i+1)%3];
  pair n1=unit(rotate(-90)*(X-Wp)),n2=unit(rotate(-90)*(Y-X));
  draw(arc(X,d,degrees(n1),degrees(n2),CCW),linewidth(0.8));
}
label("$5$",(A+B)/2,SE);label("$7$",(B+C)/2,N);label("$4$",(A+C)/2,SW);
draw((5,0)--(5,-d),Arrows(TeXHead,1.5));label("$2$ cm",(5,-d),S,fontsize(8pt));
label("\mbox{}",(min(currentpicture,user=true)+max(currentpicture,user=true))/2,black); // keeps the text after the figure black when LaTeX typesets the labels inline
\end{asy}

(The figure is not to scale; the frame is drawn much wider than $2$ cm.) The perimeter is $4+5+7=16$ m $=1600$ cm. So the frame's area is
\[
2\cdot1600+\pi\cdot2^2=3200+4\pi\approx3212.57\text{ cm}^2,
\]
which rounds to $\boxed{3213}$ square centimeters.
